%% Part II -- Basics (comic pages 11 to 19).
%% The three ingredients the book names for the current wave: computer power,
%% better algorithms, data. Then the reminder that what came out of them is
%% narrow, and what a second AI Winter would cost.
%% Fragment only: no preamble, \input by ai_for_all.tex.

\sectionsubtitle{pages 11 to 19}
\section{Basics}

% Page 13. The book refuses the word "intelligence" for what is shipping, and
% names what it would take to earn it. The footnote on the same page concedes
% it will keep saying AI anyway, which is worth reading aloud.
\begin{frame}{What is Artificial Intelligence?}
  \panel{page 13, the essence of a new situation}
  \begin{comicquote}{13}
    Up to now, it's more adequate to speak of `artificially augmented
    intelligence'.
  \end{comicquote}
  \slidesources{Schneider2019}
\end{frame}

% Page 14, first of the three key indicators. The number is the point: not
% Moore's 18 months, but 3,5, and the page prints both so the gap shows.
\begin{frame}{Computer Power}
  \panel{page 14, the cost and power curve}
  \begin{comicquote}{14}
    The amount of computer power used for training AI has gone up to 3,5
    month-doubling time.
  \end{comicquote}
  \slidesources{Schneider2019}
\end{frame}

% Page 15, second indicator. One sentence for the whole shift from rule
% writing to learning, and it is the sentence the room already half knows.
\begin{frame}{Algorithms 1}
  \panel{page 15, a cat? a cat? a dog..?}
  \begin{comicquote}{15}
    A second key indicator are better algorithms. Now we can teach machines
    instead of programming them.
  \end{comicquote}
  \slidesources{Schneider2019}
\end{frame}

% Page 16 is the buzzword page: ML, DL, NN, and the AI bookshelf that holds
% more than the three. The quote is the one that separates DL from ML, which
% is the distinction the room most often gets wrong.
\begin{frame}{Algorithms 2}
  \panel{page 16, the AI bookshelf}
  \begin{comicquote}{16}
    DL uses artificial Neural Networks (NN) that learn patterns directly
    from the data they are fed.
  \end{comicquote}
  \slidesources{Schneider2019}
\end{frame}

% Page 17, third indicator, and the widest definition of data in the book.
% The page itself is a word cloud of everything that counts as a feature.
\begin{frame}{Data 1}
  \panel{page 17, the data word cloud}
  \begin{comicquote}{17}
    Any form of raw fact or figure is data. Whether on paper or in
    electronic form.
  \end{comicquote}
  \slidesources{Schneider2019}
\end{frame}

% Page 18 adapts Maslow into a data pyramid: collect, store, transform,
% explore, AI. Four of the five steps are data work, which is the whole
% argument of the page and the reason it is quoted rather than summarised.
\begin{frame}{Data 2}
  \panel{page 18, the data pyramid}
  \begin{comicquote}{18}
    AI is a journey that begins with data: collecting, storing, transforming
    and exploring it.
  \end{comicquote}
  \slidesources{Schneider2019}
\end{frame}

% Page 19 runs the two AI Winters and lands on what today's systems actually
% are. This is the page to hold the room on: everything in Chances and Risks
% is about narrow AI, not the thing the news means.
\begin{frame}{General AI}
  \panel{page 19, the two AI Winters}
  \begin{comicquote}{19}
    Despite their impressive progress and success, today's AI is narrow. Its
    tasks are often classification and need a lot of data and a lot of
    energy.
  \end{comicquote}
  \slidesources{Schneider2019}
\end{frame}
